\documentclass[12pt,a4paper]{article}

\usepackage[a4paper,margin=1in]{geometry}
\usepackage{graphicx}
\usepackage{booktabs}
\usepackage{amsmath,amssymb}
\usepackage{float}
\usepackage{hyperref}
\usepackage{xcolor}
\usepackage{listings}
\usepackage{enumitem}
\usepackage{fancyhdr}
\usepackage{longtable}
\usepackage{mathptmx}
\usepackage[bottom]{footmisc}

\hypersetup{colorlinks=true,linkcolor=blue,urlcolor=blue}
\setlist{itemsep=2pt,topsep=3pt}

\pagestyle{fancy}
\fancyhf{}
\lhead{ICS1512}
\rhead{Machine Learning Algorithms Laboratory}
\cfoot{\thepage}

\lstset{
language=Python,
basicstyle=\ttfamily\small,
numbers=left,
frame=single,
breaklines=true,
showstringspaces=false
}

\begin{document}

\begin{tabular}{|p{4cm}|p{11cm}|}
\hline
Experiment No. & 9 \\ \hline
Experiment Title & Perceptron vs Multilayer Perceptron (A/B Experiment) with Hyperparameter Tuning\footnotemark \\ \hline
Student Name & Sanjay S \\ \hline
Register Number & 3122247001056 \\ \hline
Faculty & Dr. Poreddy Ajay Kumar Reddy \\ \hline
Submission Date & 20/09/2026 \\ \hline
\end{tabular}
\footnotetext{\textbf{GitHub Repository: }\url{https://github.com/sanjay-6727/Introdution-to-Machine-Learning-}}

\section{Objective}
Build a Perceptron (PLA) and a Multilayer Perceptron (MLP) from scratch, tune the MLP properly, and compare
the two on handwritten character recognition.

\section{Problem Statement}
Classify 3,410 handwritten character images into 62 classes (0--9, A--Z, a--z).
\begin{itemize}
\item \textbf{Model A:} single layer Perceptron, step activation.
\item \textbf{Model B:} MLP with hidden layers, trained with backpropagation.
\item Compare them on a held out test set.
\end{itemize}

\section{Dataset Description}
\begin{longtable}{|p{2.8cm}|p{3.7cm}|p{1.5cm}|p{2.6cm}|p{1.5cm}|p{1.4cm}|}
\hline
\textbf{Dataset} & \textbf{Source} & \textbf{Samples} & \textbf{Features} & \textbf{Classes} & \textbf{Missing} \\ \hline
English Handwritten Characters & Kaggle, dhruvildave\slash english-handwritten-characters-dataset & 3{,}410 & 784 (28x28 pixels) & 62 & none \\ \hline
\end{longtable}

\begin{itemize}
\item 55 images per class, so the data is balanced.
\item Split 70/15/15, stratified, seed 42: 2387 train, 511 validation, 512 test.
\item Validation set used for all tuning, test set used once at the end.
\end{itemize}

\section{Exploratory Data Analysis}
\begin{figure}[H]
\centering
\includegraphics[width=0.95\linewidth]{outputs/Experiment_9_eda_combined.png}
\caption{Sample images (top) and class distribution (bottom).}
\end{figure}

\textbf{Inference:}
\begin{itemize}
\item Every class has exactly 55 images, so chance level is only 1/62 (about 1.6\%).
\item Handwriting varies a lot inside one class (thick or thin, slanted or upright).
\item Some characters look almost the same at 28x28: 0/O/o and 1/I/l. This makes the problem hard.
\end{itemize}

\section{Data Preprocessing}
\begin{itemize}
\item Convert to grayscale, pad to a square, resize to 28x28.
\item Invert so ink = 1, then scale each image to [0, 1].
\item Flatten to 784 values.
\item Standardize each pixel using the training set mean and standard deviation.
\end{itemize}

\lstinputlisting[caption={Image preprocessing},firstline=38,lastline=43]{run_experiment.py}

\section{Machine Learning Algorithm}
Both models are written in NumPy. scikit-learn is only used for the split and the metrics.

\textbf{Model A: Perceptron (PLA)}
\begin{itemize}
\item Step activation: $\hat y = \mathrm{step}(w^Tx+b)$
\item Update rule: $w_{t+1} = w_t + \eta\,(y-\hat y)\,x$
\item 62 one vs rest perceptrons, prediction = class with the highest score.
\item Fast and simple, but only draws straight boundaries.
\end{itemize}

\textbf{Model B: MLP}
\begin{itemize}
\item Input, hidden layers with ReLU / tanh / sigmoid, softmax output.
\item Loss: cross entropy (MSE also tested). Trained with backpropagation.
\item Optimizers: Gradient Descent, SGD, Adam.
\item Learns non linear boundaries, but has many hyperparameters and overfits easily.
\end{itemize}

\lstinputlisting[caption={PLA from scratch},firstline=106,lastline=134]{run_experiment.py}
\lstinputlisting[caption={One MLP training step (backpropagation and optimizer update)},firstline=191,lastline=217]{run_experiment.py}

\section{Experimental Setup}
\begin{tabular}{ll}
\toprule
Python Version & 3.12.3 \\
NumPy Version & 2.4.4 \\
Scikit-Learn Version & 1.8.0 \\
Operating System & Ubuntu Linux, CPU only \\
Random Seed & 42 \\
Hardware & \\
\bottomrule
\end{tabular}

\section{Hyperparameter Selection}
\textbf{PLA:} the learning rate $\eta$ makes no difference. With weights starting at zero, $\eta$ only scales
the weights, and the step function and argmax ignore scale. Used $\eta = 0.01$.

\textbf{MLP:} tuned one factor at a time on validation accuracy, keeping the best value each time.
Starting point: 1 hidden layer of 128 ReLU units, cross entropy, Adam, lr $10^{-3}$, batch 64, 60 epochs.

\begin{table}[H]
\centering
\small
\begin{tabular}{lcccc}
\toprule
Factor & Settings tried & Best setting & Best val. acc. & Worst val. acc. \\
\midrule
Optimizer + learning rate & 11 & sgd lr=0.5 & 45.6\% & 15.9\% \\
Activation & 3 & relu & 45.6\% & 41.1\% \\
Cost function & 2 & cross-entropy & 45.6\% & 43.6\% \\
Batch size & 5 & bs=32 & 47.7\% & 25.6\% \\
Depth / width + lr & 10 & 1x256 & 50.1\% & 1.6\% \\
L2 weight decay & 4 & l2=0.001 & 53.6\% & 48.1\% \\
\bottomrule
\end{tabular}
\caption{Tuning summary. Detailed tables are in Section~\ref{sec:tuning}.}
\end{table}

\begin{figure}[H]
\centering
\includegraphics[width=0.95\linewidth]{outputs/Experiment_9_tuning_combined.png}
\caption{Train vs validation accuracy for every tuning stage (top), optimizer convergence and depth (bottom).}
\end{figure}

\begin{table}[H]
\centering
\small
\begin{tabular}{|p{3.2cm}|p{3.4cm}|p{8.4cm}|}
\hline
\textbf{Hyperparameter} & \textbf{Chosen} & \textbf{Why} \\ \hline
Hidden layers & 1 x 256 & Best validation accuracy (50.1\%). More layers did not clearly help. \\ \hline
Activation & ReLU & Best validation accuracy (45.6\%). \\ \hline
Cost function & Cross entropy & Better than MSE (45.6\% vs 43.6\%). \\ \hline
Optimizer, lr & SGD, 0.5 & Best of 11 combinations (45.6\%). Best Adam: 44.0\%. \\ \hline
Batch size & 32 & Best of 16 to 256 (47.7\%). \\ \hline
L2 weight decay & $10^{-3}$ & Best validation accuracy (53.6\% vs 50.1\% without). \\ \hline
Early stopping & patience 20 & Stopped after 119 epochs. \\ \hline
\end{tabular}
\caption{Final MLP configuration.}
\end{table}

\section{Results}
\begin{table}[H]
\centering
\begin{tabular}{lcccccc}
\toprule
Model & Accuracy & Precision & Recall & F1 & Macro AUC & Train acc. \\
\midrule
A: PLA & 23.0\% & 25.7\% & 23.0\% & 23.1\% & 0.782 & 76.3\% \\
B: MLP (tuned) & \textbf{49.4\%} & \textbf{50.9\%} & \textbf{49.5\%} & \textbf{48.8\%} & \textbf{0.939} & 99.7\% \\
MLP (untuned)$^\dagger$ & 1.2\% & 0.2\% & 1.2\% & 0.3\% & -- & -- \\
\bottomrule
\end{tabular}
\caption{Test set results (precision, recall and F1 are macro averages).}
\end{table}
$^\dagger$ sigmoid, MSE, SGD, lr 0.01. It stays at chance level.

\section{Result Visualization}
\begin{figure}[H]
\centering
\includegraphics[width=0.95\linewidth]{outputs/Experiment_9_results_combined.png}
\caption{Confusion matrices (PLA left, MLP right) and ROC curves on the test set.}
\end{figure}

\begin{figure}[H]
\centering
\includegraphics[width=0.85\linewidth]{outputs/Experiment_9_curves_combined.png}
\caption{Error and loss against epochs, and test metrics for both models.}
\label{fig:curves}
\end{figure}

\textbf{Inference:}
\begin{itemize}
\item PLA confusion matrix is scattered. MLP has a clear diagonal.
\item MLP mistakes are mostly look alikes: x$\to$u (3), l$\to$t (3), i$\to$1 (3), I$\to$l (3) (true $\to$ predicted, count).
\item Hardest classes: `x' (0\%), `I' (11\%), `l' (12\%), `5' (12\%). Easiest: `U' (88\%), `c' (78\%), `C' (75\%), `3' (75\%).
\item ROC: MLP is close to the top left corner (AUC 0.939), PLA is much lower (0.782).
\item PLA training error stops at about 24\% and never reaches zero. MLP training loss goes almost to zero.
\end{itemize}

\section{Detailed Tuning Results}\label{sec:tuning}
Each run used 60 epochs. ``--'' means the run never reached exactly zero training error.

\begin{table}[H]\centering\small
\begin{tabular}{lcccc}\toprule
Setting & Train acc. & Val. acc. & Epochs to 0 train err. & Final train loss \\\midrule
gd lr=0.1 & 24.0\% & 15.9\% & -- & 3.4027 \\
gd lr=0.5 & 60.6\% & 36.2\% & -- & 1.7591 \\
gd lr=1 & 76.8\% & 40.3\% & -- & 1.1169 \\
sgd lr=0.01 & 54.7\% & 32.3\% & -- & 2.0433 \\
sgd lr=0.05 & 89.2\% & 43.8\% & -- & 0.5896 \\
sgd lr=0.1 & 97.7\% & 44.8\% & -- & 0.2072 \\
sgd lr=0.5 & 99.5\% & \textbf{45.6\%} & -- & 0.0233 \\
adam lr=0.0001 & 61.3\% & 35.2\% & -- & 1.7563 \\
adam lr=0.001 & 99.6\% & 44.0\% & -- & 0.0806 \\
adam lr=0.003 & 100.0\% & 44.0\% & 46 & 0.0070 \\
adam lr=0.01 & 100.0\% & 42.5\% & 34 & 0.0013 \\
\bottomrule\end{tabular}
\caption{Stage 1 -- optimiser and learning rate (60 epochs; ``--'' = never reached zero training error).}\label{tab:opt}\end{table}
\begin{table}[H]\centering\small
\begin{tabular}{lcccc}\toprule
Setting & Train acc. & Val. acc. & Epochs to 0 train err. & Final train loss \\\midrule
relu & 99.5\% & \textbf{45.6\%} & -- & 0.0233 \\
tanh & 99.9\% & 42.1\% & 46 & 0.0339 \\
sigmoid & 90.6\% & 41.1\% & -- & 0.5108 \\
\bottomrule\end{tabular}
\caption{Stage 2 -- activation function.}\label{tab:act}\end{table}
\begin{table}[H]\centering\small
\begin{tabular}{lcccc}\toprule
Setting & Train acc. & Val. acc. & Epochs to 0 train err. & Final train loss \\\midrule
cross-entropy & 99.5\% & \textbf{45.6\%} & -- & 0.0233 \\
MSE & 85.3\% & 43.6\% & -- & 0.1864 \\
\bottomrule\end{tabular}
\caption{Stage 3 -- cost function.}\label{tab:loss}\end{table}
\begin{table}[H]\centering\small
\begin{tabular}{lcccc}\toprule
Setting & Train acc. & Val. acc. & Epochs to 0 train err. & Final train loss \\\midrule
bs=16 & 56.5\% & 25.6\% & -- & 8.8025 \\
bs=32 & 100.0\% & \textbf{47.7\%} & 38 & 0.0045 \\
bs=64 & 99.5\% & 45.6\% & -- & 0.0233 \\
bs=128 & 99.6\% & 46.6\% & -- & 0.0396 \\
bs=256 & 98.5\% & 45.0\% & -- & 0.1326 \\
\bottomrule\end{tabular}
\caption{Stage 4 -- batch size.}\label{tab:batch}\end{table}
\begin{table}[H]\centering\small
\begin{tabular}{lcccc}\toprule
Setting & Train acc. & Val. acc. & Epochs to 0 train err. & Final train loss \\\midrule
1x128 lr=0.5 & 100.0\% & 47.7\% & 38 & 0.0045 \\
1x128 lr=0.1 & 99.5\% & 45.4\% & -- & 0.0597 \\
2x128 lr=0.5 & 57.0\% & 24.7\% & -- & 1.9932 \\
2x128 lr=0.1 & 100.0\% & 46.0\% & 57 & 0.0099 \\
3x128 lr=0.5 & 1.6\% & 1.6\% & -- & 4.1272 \\
3x128 lr=0.1 & 100.0\% & 48.7\% & 45 & 0.0034 \\
4x128 lr=0.5 & 33.5\% & 18.0\% & -- & 2.9929 \\
4x128 lr=0.1 & 100.0\% & 49.5\% & 34 & 0.0016 \\
1x256 & 99.9\% & \textbf{50.1\%} & 28 & 0.0057 \\
1x512 & 100.0\% & 49.7\% & 27 & 0.0038 \\
\bottomrule\end{tabular}
\caption{Stage 5 -- depth/width and learning rate.}\label{tab:depth}\end{table}
\begin{table}[H]\centering\small
\begin{tabular}{lcccc}\toprule
Setting & Train acc. & Val. acc. & Epochs to 0 train err. & Final train loss \\\midrule
l2=0 & 99.9\% & 50.1\% & 28 & 0.0057 \\
l2=0.0001 & 99.8\% & 49.7\% & 44 & 0.0130 \\
l2=0.001 & 99.5\% & \textbf{53.6\%} & -- & 0.0842 \\
l2=0.01 & 73.5\% & 48.1\% & -- & 1.1224 \\
\bottomrule\end{tabular}
\caption{Stage 6 -- L2 weight decay.}\label{tab:l2}\end{table}

\section{Discussion}
\textbf{Q1. Why does PLA underperform MLP?}
\begin{itemize}
\item It draws one straight line per class, but these classes are not linearly separable.
\item It cannot even fit its own training data (76.3\% train accuracy).
\item The step update has no loss or margin, so on this data it never settles.
\item Result: 23.0\% test accuracy vs 49.4\% for the MLP.
\end{itemize}

\textbf{Q2. Which hyperparameters mattered most?}
\begin{itemize}
\item \textbf{Learning rate} by far. Validation accuracy ranged from 15.9\%--45.6\% across optimizer and lr settings.
\item lr 0.5 broke batch size 16 (25.6\%) and 3 hidden layers (1.6\%).
\item L2 weight decay helped by 3.5 points.
\item Activation and loss: small differences, within noise (about $\pm$2 points).
\end{itemize}

\textbf{Q3. Did SGD vs Adam affect convergence?}
\begin{itemize}
\item Adam worked well over a wide lr range ($10^{-3}$ to $10^{-2}$, validation 42.5\%--44.0\%).
\item SGD was very sensitive: lr 0.01 gave 32.3\%, lr 0.5 gave 45.6\%.
\item Full batch GD was slowest (one update per epoch).
\item Best SGD and best Adam differ by only 1.6 points, so neither is clearly better.
\end{itemize}

\textbf{Q4. Did more hidden layers always help?}
\begin{itemize}
\item No. At lr 0.5, 3 layers failed completely (1.6\%).
\item At lr 0.1, going from 1 to 4 layers gave about 4 points (45.4\% to 49.5\%), close to noise.
\item One wide layer (50.1\%) did as well, so I kept one hidden layer.
\item Every stable model already reached about 100\% train accuracy, so extra depth had nothing left to fix.
\end{itemize}

\textbf{Q5. Did the MLP overfit? How to reduce it?}
\begin{itemize}
\item Yes. Train accuracy 99.7\%, test accuracy 49.4\%.
\item Without regularization, validation loss is lowest at epoch 8, then rises while training loss keeps falling.
\item Used: L2 weight decay and early stopping.
\item Could also try: data augmentation, dropout, more data, or a CNN.
\end{itemize}

\textbf{Impact of tuning:} untuned MLP 1.2\%, tuned MLP 49.4\%. During the search, validation accuracy went from
44.0\% to 53.6\%. The final model scores 56.0\% on validation but 49.4\% on test, because the settings were
picked on the validation set. The test number is the honest one.

\section{Conclusion}
\begin{itemize}
\item The PLA reaches 23.0\%. The tuned MLP reaches 49.4\% (macro F1 48.8\%).
\item Learning rate was the most important hyperparameter. Deeper was not clearly better.
\item Overfitting is the main problem left, because there are only about 38 training images per class.
\end{itemize}

\section{Learning Outcomes}
\begin{itemize}
\item Implemented the PLA and an MLP with backpropagation from scratch.
\item Learned to tune hyperparameters on a separate validation set.
\item Saw how learning rate, batch size and depth affect each other.
\item Learned to read confusion matrices, ROC curves and convergence curves.
\end{itemize}

\section{References}
\begin{enumerate}
\item Stephen Marsland, ``Machine Learning - An Algorithmic Perspective'', 2nd Edition, Chapman and
Hall/CRC Machine Learning and Pattern Recognition Series, 2015.
\item Ethem Alpaydin, ``Introduction to Machine Learning'', 3rd Edition, The MIT Press, 2014.
\item F. Rosenblatt, ``The perceptron: a probabilistic model for information storage and organization in the brain'',
Psychological Review, 65(6), 1958.
\item D. P. Kingma and J. Ba, ``Adam: a method for stochastic optimization'', ICLR 2015.
\end{enumerate}

\end{document}
